\chapter{C14.1b: the bounded testing-ratio complement at every unit $h$}
\label{component:p5cd5fa89abcc}
\noindent\textbf{Component }\path{p5cd5fa89abcc}\par
{\small\noindent Source: \path{BASE/proofs/c1_unbounded/c141b_bounded_testing_ratio.tex}\par}
{\small\noindent SHA256: \path{5cd5fa89abcc09162c1c8ff3f19cc62329522b5890615df07ba5de404204dccd}\par}
{\small\noindent Inventory use: BC1-BC11 exact bounded state/testing-ratio admission and retained complement | BC1--BC11 | C14\_\allowbreak{}ENTIRE\_\allowbreak{}ORIGINAL\_\allowbreak{}PARENT\_\allowbreak{}GRAPH; precise subordinate uses retained below\par}
\medskip
\section{C14.1b: the bounded testing-ratio complement at every unit $h$}
\label{p5cd5fa89abcc:sec:c141b-bounded-testing-ratio}

This is a further partial-domain admission, independent of the positive
angular-pivot restriction in the preceding state-free proof. Keep the
entire original physical C14.1b source, all strict gaps, original scale
and prior-complement intersections, and all allowed coefficient openings.
Thus $x,z=1/L,\sigma=1/\eta$ are positive infinitesimals, $H=h/\sigma$
is unbounded, and all source quantities belong to $K_c$, independently
of tests and states in $K_t$. We assume here that $h$ is a positive
bounded unit. The claim is about the complete original P4/SP6 parent.

After the inherited literal-axis, state-sign and testing-sign exits,
consider $u,w,W,Z>0$, with $u,w$ infinitesimal, and define exclusively
in the testing field
\[
                 c=\frac{uW}{wZ}>0.
                 \tag{BC1}
\]
The new assertion is that every real state and test with $c$ bounded
is admitted. In particular this includes infinitesimal $c$, with no
positive real lower bound on $c$. In the actual retained C14.1b
complement one can additionally keep $W/Z$ unbounded; nothing below
needs that condition. A complex state is covered whenever each of its
real and imaginary parts belongs to this domain or to an already
admitted domain. No condition is imposed on a nonexistent complex
order or complex state ratio.

\paragraph{Exact bounded source input.}
Use SF1--SF7 and the complete radial extraction SF8. They were proved
with $z$ infinitesimal and do not require a positive lower bound on $z$.
In particular
\[
 \lambda>0,\quad \lambda\sigma\sim2h,\quad
 \rho=\lambda x/\sigma\text{ bounded},\quad
 \lambda x\text{ infinitesimal},\quad K_A\text{ bounded}.
 \tag{BC2}
\]
No positivity of $K_A$ is assumed. The exact sphere split gives
a positive original Gamma pivot when $d_0$ has positive residue.
When $d_0$ is infinitesimal, $\rho$ is a positive unit and
$\sigma/k$ is bounded, where $k=\sqrt x$. The actual second
Gamma normal has the source factorization SF18--SF20, with
$R_0/\sigma$ bounded. These facts include literal $D=0$.

\paragraph{The complete scalar circuit and a new radial normalization.}
Write $q=\sqrt z$ and use exactly the three nonnegative separated
monomials from SF10:
\[
 a_1=\sqrt{\lambda\sigma}\,uZ,\qquad
 a_2=\sqrt\rho\,uwW,\qquad
 a_3=\sqrt\lambda\,xqW,
 \qquad G=(a_1a_2a_3)^{1/3},\quad d_i=a_i-G.
 \tag{BC3}
\]
Every fractional power is the product of its pure-source and
pure-testing powers. The exact CT4 scalar identity applies without
changing the source or ordering a mixed ratio: the scalar packet
$\mathscr P$ of SF11 plus $\lambda xR^2$ is
$\sum_i a_i^2-3G^2$ plus the two printed nonnegative CT4 squares.
For $p_i=a_i^{1/3}$, its exact strict circuit decomposition is
\[
 \sum_i a_i^2-3G^2
 =\sum_i d_i^2+
   \sum_{j=1}^3\sum_{i<k}
       \{\sqrt{Gp_j}(p_i-p_k)\}^2.
 \tag{BC4}
\]
The nine supplementary squares and the two CT4 squares remain whole.
Thus the circuit supplies three unit diagonal pivots on $d_1,d_2,d_3$;
its only cost is the infinitesimal source correction $-\lambda xR^2$.

Put
\[
 k_c=z^{1/3}c^{1/3},\qquad
 \overline J=(k_c-1)\frac{wZ}{q}.
 \tag{BC5}
\]
The bar in $\overline J$ is part of the root's name, not a residue.
It is a finite separated expression. Its exact relation to the
original radial root is
\[
 R=\frac{wZ}{q}-uqW
       =-(1+k_c+k_c^2)\overline J.
 \tag{BC6}
\]
No mixed expression on the right is inverted. Because $c$ is bounded,
$k_c$ and $k_c^2$ are bounded separated coefficients with zero real
residue. Therefore the coefficient in BC6 has residue $-1$.
After BC6 the clipped SF8 radial form is a bounded strict Gram on
$(N_y,\overline J)$, including the $b$ and $Q$ opening endpoints
in residue. All exact positive radial excesses are retained.

This is precisely the rescaling
$\overline J=J_c/(q c^{2/3})$ of CT6, but definition BC5 avoids
an inverse of $c$ even in the displayed root. It is valid when
$c$ is infinitesimal.

\paragraph{The signed angular packet has a small exact conductor.}
The following whole-state identity is the rescaled CT13 conductor:
\[
\begin{split}
 A_\theta={}&
 \frac{qcw}{\sqrt{\lambda\sigma}}d_1
 +\frac{k_c^2}{\sqrt\lambda}d_2
 -\frac1{\sqrt\lambda}d_3\\
 &+(k_c+1)\sqrt x\,z^{2/3}c^{2/3}\overline J
 =:\mathcal V_A.
\end{split}\tag{BC7}
\]
All four coefficients are finite sums of bounded separated factors
with zero residue. The first uses bounded $c$, infinitesimal $q,w$
and the positive source unit $\lambda\sigma$. The second and third
use $1/\sqrt\lambda$ infinitesimal. For the last, expand it as
$\sqrt x\,z^{2/3}c^{2/3}+\sqrt x\,zc$.
No bound involving an ordered product of independent rates occurs.

One can therefore replace the possibly signed angular coefficient
by a strict reserve, using an identity rather than an inequality.
For the SF13 root $T=kq(u-k)W$, exactly
\[
\begin{split}
 K_AA_\theta^2={}&\tfrac12 A_\theta^2+\tfrac12 T^2
       +\tfrac12(A_\theta^2-T^2)
       +(K_A-1)\mathcal V_A^2,\\
 A_\theta^2-T^2={}&q^2(u-k)^2W^2(u^2+2uk).
\end{split}\tag{BC8}
\]
The last expression in the second line is an exact separated SOS.
The first two rows in the first line give strict pivots $1/2$ on
$A_\theta,T$. The signed final term has bounded infinitesimal Gram
entries on $(d_1,d_2,d_3,\overline J)$ by BC7 and bounded $K_A$.
In particular this replacement remains valid at $\operatorname{res}h=1/2$
and when $K_A$ has negative residue. It is paid exactly once.

\paragraph{Every normal of the original parent.}
Keep the original $N_x$ pivot and the exact SF8 radial extraction.
If $d_0$ has positive residue, retain the strict original Gamma
pivot and use its two full normals SF16 directly. Their coefficients
are bounded infinitesimals, including after BC6. If $d_0$ is
infinitesimal, use SF15 and the complete second-normal factorization
SF20, just as in the state-free proof. Their circuit roots are now
\[
 D_1=(d_1-d_2)/\sqrt2,\qquad
 D_2=(d_2-d_3)/\sqrt2.
 \tag{BC9}
\]
Thus no new division by a sphere coordinate is needed. SF7 and SF19
make every resulting coefficient bounded infinitesimal; $A_\theta,T$
have the strict reserves just supplied by BC8. The positive original
$4D\Gamma^2/\sigma^2$ row is retained whole in this latter case.

All other SP6 terms are precisely the five complete crosses in SF9,
the infinitesimal radial remainder CT7a, and the cost
$-\lambda xR^2$. Substitute BC6 where $R$ occurs. They remain
bounded infinitesimal perturbations of the strict frame
\[
 (N_x,N_y,\overline J,A_\theta,T,d_1,d_2,d_3),
 \tag{BC10}
\]
with $\Gamma$ adjoined in the positive-$d_0$ case. The source-unit
radial block, fixed $N_x$ pivot, BC8 angular pivots, and BC4 circuit
pivots give a positive-definite real residue. Retain exactly once
$4\beta D_wN_x^2$, $2(2D+1)(uX_e)^2$, $4bpzu^2\Gamma^2$,
all nonconstant positive radial packets and source-weight excesses,
the nine BC4 squares, the two CT4 squares, and the positive BC8
domination remainder. No full state normal has been discarded.
The separated bounded-Gram lemma now gives the complete parent
as tensor SOS on the asserted domain.

The proof concerns every independent test/state in BC1 with $c$
bounded. Literal axes and opposite-sign configurations keep their
inherited admissions. Applying the real proof or an earlier exit
separately to the real and imaginary parts gives the complex-state
join stated at the start; it does not enlarge the real domain by
assuming an order on complex amplitudes.

\paragraph{Exact improvement and retained obstruction.}
Combining this result with the previously accepted source subdomain
$\operatorname{res}h<1/2$, the real C14.1b remainder can be narrowed
to its original retained source/test intersections with
\[
 \operatorname{res}h\ge1/2,\qquad
 u,w>0\text{ infinitesimal},\quad WZ>0,\quad W/Z\text{ unbounded},
 \quad \frac{uW}{wZ}\text{ unbounded}.
 \tag{BC11}
\]
This complement still contains the difficult source-contact scale:
on the actual source contact $uW/(wZ)=L=1/z$, which is unbounded.
For unbounded pure-testing $c$, BC7's coefficients involving $c$
cannot be declared bounded separated coefficients, and the present
argument does not cover that case. No full C14.1b closure or
nonemptiness assertion is made.

The narrow checker verifies the exact radial relation, the complete
angular conductor, the circuit decomposition and conversion, and
the once-paid angular reserve. The source bounds, every remaining
normal and the independent-field Gram argument are the analytical
proof above, with precisely the SF and CT dependencies identified.
